\section{What a strict count comparison actually needs}
For a centered variance-one law put $p=\Pp(X>0)$, $q=1-p$,
$m=\E X_+$, $\ell=m/(pq)$, and
$X=\ell(B+U-p)$, where $B=\mathbf1_{X>0}$ and
$\E(U\mid B)=0$. Write $s=\E U^2$ and $\lambda=ns$.
Let $np+H/\ell=k+u$, with an interior integer count $k$ and
$|u|\le1/2$. All variables are constructed from the same marginal.

For fixed target $c=\CE$, define
$a=\sqrt n\{F_n(H)-\Phi(H/\sqrt n)\}-c\rho$.
The exact count identity has left side
\[
                     a+\rho_B(c-C_B),
\]
where $C_B$ is the actual matched-binomial score. If $a\ge0$ and
the external Bernoulli theorem gives $C_B\le c$, this is nonnegative.
Thus a whole-count upper bound of the form
\[
 \sqrt n\pi_k\lambda B(\lambda)+\frac{\lambda}{n}I(\lambda)
\]
is contradictory for $\lambda>0$ whenever valid constants satisfy
\begin{equation}\label{ex:final}
 B(\lambda)<0,\qquad
 m_0B(\lambda)+\frac{\max(0,I(\lambda))}{N}<0,
 \quad \sqrt n\pi_k\ge m_0>0,
 \quad n\ge N.
\end{equation}
The first sign gate is necessary before applying a lower mass bound.
Zero noise is handled separately by the Bernoulli theorem.

Neither $D<0.784$ nor a raw Fourier bound below $2.24$ occurs in
\eqref{ex:final}. Those are prescribed outputs of the retained initial
acceptor. They supply later acquisition inputs through its particular
dependency chain; they are not necessary scalar constants in every possible
count proof. Conversely, obtaining a smaller $D$ alone does not establish
the conditional support, amplitude, marked-crossing and Gaussian-repair
hypotheses needed to produce $B,I$. The example below shows why.

There is a useful exact test for the remaining marked parameter. On a
fixed positive noise value, the retained whole-count formula has the form
\[
 B(\lambda)=z+\frac{A_\lambda}{z}-F_\lambda,
 \quad A_\lambda\ge0,\quad0<z\le1/U,
\]
where $\eta=z\lambda\le1$ on the noise band $[L,U]$.
Write $F=F_\lambda-\max(0,I(\lambda))/(Nm_0)$ and $Z=1/U$.
Then \eqref{ex:final} follows at this noise value precisely when
$z+A_\lambda/z<F$ for an allowed $z$.
For $A_\lambda>0$ the minimum is
\begin{equation}\label{ex:onevariable}
 \min_{0<z\le Z}\left(z+\frac{A_\lambda}{z}\right)
 =\begin{cases}
    2\sqrt{A_\lambda},&Z^2\ge A_\lambda,\\
    Z+A_\lambda/Z,&Z^2<A_\lambda.
   \end{cases}
\end{equation}
This follows by differentiating, or completing
$(\sqrt z-\sqrt{A_\lambda/z})^2$ and checking the boundary.
If $A_\lambda=0$, feasibility is exactly $F>0$; the infimum zero is
not attained. In the interior case the exact rational gate is
$F>0$ and $F^2>4A_\lambda$. Equality does not suffice.
For a full closed noise band a common $z$ must still satisfy both
endpoint gates before using the existing convexity proof; optimizing $z$
separately at the endpoints does not establish that common witness.
In particular a zero fractional-part gain makes this whole-count formula
nonnegative for every $z>0$, even if the mark cost vanishes. It cannot be
repaired simply by shrinking the weighted resource.

\section{The harmonic argument works up to span one}
For a real number $x$, write $\{x\}=x-\lfloor x\rfloor$.

\begin{lemma}[Extended conditional-span domain]\label{ex:harmonic}
Let $Z_1,\ldots,Z_j$ be independent real variables, each supported in an
interval of length at most $b$, where $0<b<1$.
Put $W=\sum_i Z_i$ and $t=\sum_i\operatorname{Var}(Z_i)$.
Suppose $K>0$ satisfies
\[
                  1-\cos(2\pi x)\le Kx\quad(0\le x\le1).
\]
For every $w\in\R$ and every
$0<r\le2\sin^2(\pi b)/b^2$,
\begin{equation}\label{ex:gain}
                     \E\{w-W\}\ge K^{-1}(1-e^{-rt}).
\end{equation}
In particular $K=8$ is valid without a finite scalar certificate.
The inherited value $K=23/5$, or its accepted improvement, may be retained
when its separate scalar majorization has been proved.
\end{lemma}
\begin{proof}
Let $Z_i'$ be an independent copy. Since $|Z_i-Z_i'|\le b$ and sine is
concave on $[0,\pi]$, the chord from $0$ to $\pi b$ gives
\[
 \sin^2\bigl(\pi(Z_i-Z_i')\bigr)
 \ge\frac{\sin^2(\pi b)}{b^2}(Z_i-Z_i')^2.
\]
Consequently, writing $v_i=\operatorname{Var}(Z_i)$,
\[
 |\E e^{2\pi iZ_i}|^2
 =1-2\E\sin^2\bigl(\pi(Z_i-Z_i')\bigr)
 \le1-\frac{4\sin^2(\pi b)}{b^2}v_i
 \le e^{-2rv_i}.
\]
The intermediate right side is nonnegative because it bounds the
nonnegative left side. Taking square roots and multiplying, using
independence, gives $|\E e^{2\pi iW}|\le e^{-rt}$. Periodicity then gives
\[
 K\E\{w-W\}\ge1-\Re\E e^{2\pi i(w-W)}
                \ge1-|\E e^{2\pi iW}|\ge1-e^{-rt}.
\]
No centering or continuous distribution assumption was used. To verify
$K=8$, the function $y-1+\cos y$ is nondecreasing for $y\ge0$, since its
derivative is $1-\sin y\ge0$. Thus $1-\cos(2\pi x)\le2\pi x<8x$ for
$x>0$; $\pi<4$ also follows directly from
$\pi=4\int_0^1(1+y^2)^{-1}\,dy<4$. Equality at $x=0$ is harmless.
\end{proof}

The earlier count proof restricts $b<1/2$ at this step. The same proof in
fact gives Lemma~\ref{ex:harmonic} throughout $b<1$. For example,
$b=3/4$ permits the exact rate $r=16/9$.
This extends only the harmonic component. It does not extend other
denominator, binomial, support or noise gates. In the retained whole-count
notation a sufficient acquired span is $b=A+d_0$; it must be a bound on
each actual conditional residual, not on its variance.

For example, on a phase interval $[u_l,u_h]$, the retained individual
first-mark denominators are positive when
$1-2d_0-u_h>0$ and $1-A+u_l>0$, with $c=1-d_0$.
On the full phase interval $[-1/2,1/2]$, these become $d_0<1/4$
and $A<1/2$. They imply $A+d_0<3/4<1$, so the extended harmonic
lemma adds no further span restriction there. The optional common-mark
bound has additional gates. None of these inequalities alone establishes
the final score, and a larger span gives a smaller harmonic rate that must
be used in that score.

\begin{proposition}[The boundary is a real obstruction]\label{ex:boundary}
No strictly positive uniform rate in \eqref{ex:gain} is possible for all
independent positive-variance summands of support span at most one.
Even a uniformly positive rate as $b\uparrow1$ is impossible.
\end{proposition}
\begin{proof}
For $0<\varepsilon<1$, let $Z$ equal $\varepsilon-1$ with probability
$\varepsilon$, and $\varepsilon$ otherwise. It is centered, has positive
variance $\varepsilon(1-\varepsilon)$ and span one.
For $j$ independent copies, $W=j\varepsilon-M$ with integer-valued
$M$, so $\{j\varepsilon-W\}=0$ identically. The right side of
\eqref{ex:gain} is positive for every $r>0$ and finite $K>0$.

For the boundary limit, the centered variable taking values $\pm b/2$
with equal probabilities has variance $b^2/4$ and characteristic modulus
$|\cos(\pi b)|$. As $b\uparrow1$, this modulus tends to one while the
variance tends to $1/4$. A positive uniform exponential decay rate would
contradict these limits. Directly for \eqref{ex:gain}, choose $w=-b/2$:
then $\E\{w-Z\}=(1-b)/2\to0$, also contradicting a positive uniform
rate with finite fixed $K$. In fact $1-|\cos(\pi b)|$ vanishes quadratically
in $1-b$, consistently with the rate in Lemma~\ref{ex:harmonic}.
\end{proof}

\section{An actual iid law with arbitrarily small moment resources}
The obstruction can be realized by the actual conditional laws of a
centered, variance-one iid marginal with a large central gap. It is not
merely an optimizer of a moment relaxation.

For rational $\varepsilon,t>0$, set
\begin{equation}\label{ex:family}
 q=\frac{1+\varepsilon-\varepsilon^2}{1+t^2},\quad
 p=1-q,\quad r=tq,
 \qquad 0<\varepsilon<p<1.
\end{equation}
Let $Y$ take the values $\varepsilon-1,\varepsilon,1$ with respective
probabilities $q\varepsilon,q(1-\varepsilon),p$, and put $X=(Y-p)/r$.
Direct calculation gives
\[
 \E Y=p,\quad
 \operatorname{Var}(Y)=q\{p+\varepsilon(1-\varepsilon)\}=r^2.
\]
Therefore $X$ is exactly centered and standardized. The first two atoms
are negative, and the third is positive. In the actual sign split,
$\ell=1/r$, $U_1=0$, and
\[
 U_0=\begin{cases}\varepsilon-1,&\text{probability }\varepsilon,\\
                  \varepsilon,&\text{probability }1-\varepsilon.
       \end{cases}
\]
Thus the actual conditional amplitude is $1-\varepsilon$, its span is
one, the normalized central-gap increment is $c_{\rm gap}=1-\varepsilon$,
and $\lambda=nq\varepsilon(1-\varepsilon)$.
In particular the central gap is wide even when the harmonic estimate
fails completely.

For a fixed interior count $k$, the actual conditional sum has the law
\[
 V_k=(n-k)\varepsilon-\operatorname{Bin}(n-k,\varepsilon).
\]
With $w=(n-k)\varepsilon$, its fractional-part gain is zero. For the four
examples checked below, $n=42$, $k=17$ and
$u=\eta=w/2$ give a legal phase $|u|\le1/2$, $0<\eta\le1$, and
$w=u+\eta$, exactly as in the marked-crossing argument. Its own-count
variance is $(n-k)\varepsilon(1-\varepsilon)>0$.

The physical defect is computed from the actual gaps, not from the
two-point skeleton:
\[
 \rho=\E|X|^3,\quad\gamma=\E X^3,\quad
 J=\sum_{x_i<x_{i+1}}
       \left(\sum_{x_j\ge x_{i+1}}p_jx_j\right)(x_{i+1}-x_i)^2,
 \qquad D=(3+\sqrt{10})\rho-\gamma-3J.
\]
The exact checker uses rational moments and the directed enclosure
\[
 3.162277660168379331<\sqrt{10}<3.162277660168379332,
\]
verified by squaring positive rationals. Its first law, with
$\varepsilon=1/100$, $t=21/25$, is explicitly
\[
\begin{array}{c|ccc}
x&-198687/70693&-169661/212079&25/21\\\hline
\Pp(X=x)&10099/1705600&999801/1705600&6957/17056.
\end{array}
\]
It has $|X|<2.811$, $1.11975<\rho<1.11976$, and
$0.67024<D<0.670241$. The four fixed examples satisfy:
\begin{center}\small
\begin{tabular}{@{}cccc@{}}\toprule
$\varepsilon$ & $t$ & strict upper bound for $D$ & strict upper bound for $\lambda$\\\midrule
$10^{-2}$ & $21/25$ & $0.670241$ & $0.246199$\\
$10^{-3}$ & $21/25$ & $0.070624$ & $0.024625$\\
$10^{-6}$ & $21/25$ & $0.000738$ & $0.000025$\\
$10^{-9}$ & $849/1000$ & $0.0000001$ & $0.0000001$\\\bottomrule
\end{tabular}
\end{center}
Every example satisfies its own recorded v6 moment-cell defect bound,
weighted $nD$ bound, and strict actual-gap criterion, as well as $D<0.784$.
This is a fresh exact comparison with the pinned table, not a replay of
the Fourier calculation that originally produced that table.

The last example also meets every listed moment, probability and inner-gap
bound of the retained response theorem's central row 19. Its sole failed
listed domain bound is the support radius $1.84$: its rare outer atom has
magnitude greater than $2.875$. Thus that theorem cannot be applied on
the strength of its favorable moment bounds alone.

\begin{corollary}[No moment-only amplitude implication]\label{ex:arbitrary}
For every $\delta>0$ there is an actual centered variance-one finite law
supported in $[-96/25,96/25]$, with positive conditional noise,
$D<\delta$ and $\lambda<\delta$ at $n=42$, whose negative conditional
residual span is exactly one and whose fractional-part gain vanishes
at a legal marked phase. The normalized central-gap increment can be
arbitrarily close to one.
\end{corollary}
\begin{proof}
At $\varepsilon=0$ the positive-weight part of \eqref{ex:family} is
the standardized two-point law at $-t,1/t$, with positive probability
$p=t^2/(1+t^2)$. Its defect is
\[
 D_B=\frac{(6+2\sqrt{10})(p-p_*)^2}{\sqrt{p(1-p)}},
 \qquad p_*=(4-\sqrt{10})/2.
\]
This follows by substituting $\rho_B,\gamma_B$ and $J_B=t+1/t$.
Choose positive rational $t$ arbitrarily near
$\sqrt{p_*/(1-p_*)}$. Then $D_B$ is arbitrarily small.
For this fixed $t$, let positive rational $\varepsilon$ tend to zero.
All moments are continuous. The extra outer-gap weight tends to zero,
so the displayed finite gap formula gives $D\to D_B$ as well;
$\lambda\to0$ while the conditional span remains one.
At the limiting value of $t$, the outer atom tends to
$-(1+2t^2)/t$. Here $0.84<t<0.85$, so its magnitude is less than
$2(0.85)+1/0.84<3$. The support is therefore within
$[-96/25,96/25]$ in a neighborhood. The sign, phase and positive-noise
conditions also hold for all sufficiently small $\varepsilon$.
\end{proof}

These laws are not asserted to satisfy a maximizing contact inequality,
nonnegative fixed-target excess, or all new signed-acquisition outputs.
They do not refute the sharp conjecture or the retained local theorem.
They refute the proposed implication from small moment resources and an
actual central gap alone to a positive fractional-part gain or a small
conditional amplitude. Contact-based confinement or a quantitatively
proved replacement for harmonic nonresonance remains an essential missing
edge in that shortcut.
